\section{Notations and basic facts} \label{section2}

The purpose of this section is to establish notations which will be
used in the sequel and recall some basic def\/initions and results
which we shall freely use afterwards without mentioning them
explicitly.

 Let $G$  be a connected and simply
connected nilpotent Lie group with Lie algebra $\frak g$, then the
exponential map ${\rm exp}: \frak g \rightarrow G$ is a
dif\/feomorphism. Let $\log: G \rightarrow \frak g$
denote the inverse of~${\rm exp}$.

\subsection{Induced representation}\label{section2.1} %\label{induced-representation}

Starting from a closed
subgroup $H$ of a nilpotent Lie group $G$ and a unitary
representation $\sigma$ of~$H$ in a Hilbert space~$\Hi_\sigma$, let
us construct a unitary representation of~$G$. We realize the unitary
representation $\ind_H^G\sigma$ of~$G$, the representation induced
on~$G$ from the representation $\sigma$ of~$H$, by left translations
on the completion of the space of continuous functions~$F$ on~$G$
with values in~$\Hi_\sigma$ satisfying
\begin{gather}\label{covarience}
F(gh) = \sigma(h^{-1})(F(g)),\qquad  g\in G, \quad h\in H,
\end{gather}
and having a compact support modulo $H$, provided with the norm
\begin{gather*}%\label{norm}
\|F\| = \left(\nu_{G, H}\big(\|F\|^2\big)\right)^{\frac{1}{2}} =
\left(\int_{G/H} \|F(g)\|^2 d\nu_{G, H}(g)\right)^{\frac{1}{2}},
\end{gather*}
where
\begin{gather*}
\big(\big(\ind_H^G\sigma\big)(g) F\big)(x) = F\big(g^{-1} x\big),\qquad g, x\in G.
\end{gather*}

\subsection{The orbit theory}\label{section2.2}

Suppose $G$ is a nilpotent Lie group with Lie algebra $\frak g$. $G$~acts on $\frak g$ (respectively $\frak g^*$) by the adjoint
(respectively co-adjoint) action. For $l\in \frak g^*$, let
\[
\frak
g(l) = \{X\in \frak g:\; \langle l, [X, \frak g]\rangle = \{0\}\}
\]
be the stabilizer of $l$ in $\frak g$ which is actually the Lie
algebra of the Lie subgroup
\[
G(l) = \{g\in G:\ g\cdot l = l\}.
\]
So, it is clear that $\frak g(l)$ is the radical of the
skew-symmetric bilinear form $B_l$ on $\frak g$ def\/ined by
\begin{gather*}%\label{skew-symmetric bilinear}
B_l(X, Y) = \langle l, [X, Y]\rangle,\qquad X, Y \in \frak g.
\end{gather*}
A subspace $\frak b(l)$ of the Lie algebra $\frak g$ is called a
polarization for $l\in \frak g^*$ if it is a maximal dimensional
isotropic subalgebra with respect to $B_l$. We have the following
equality
\begin{gather}\label{dim-pilarisation}
\dim{\frak b(l)}=\dfrac{1}{2} \big(\dim{\frak g}+\dim{\frak
g(l)}\big).
\end{gather}

If $\frak b(l)$ is a polarization for $l\in \frak g^*$ let
$B(l)=\exp(\frak b(l))$ be the connected subgroup of $G$ with Lie
algebra $\frak b(l)$ and def\/ine a character of $B(l)$ by the formula
\begin{gather*}%\label{Defn-charcter}
\chi_l(\exp(X))= e^{2i \pi \langle l, X\rangle},\qquad \forall\, X\in \frak b(l).
\end{gather*}



Now, we recall the Kirillov orbital parameters. We denote by $\hat
G$ the unitary dual of $G$, i.e.\ the set of all equivalence classes
of irreducible unitary representations of $G$. We shall sometimes
identify the equivalence class $[\pi]$ with its representative $\pi$
and we denote the equivalence relation between two representations~$\pi_1$ and $\pi_2$ by $\pi_1 \simeq \pi_2$ or even by $\pi_1 =
\pi_2$.  The dual space $\hat G$ of $G$ is parameterized canonically
by the orbital space $\frak g^*/G$. More precisely, for~$l \in \frak
g^*$ we may f\/ind a real polarization $\frak b(l)$ for~$l$.
 Then the representation $\pi_l =
\ind_{B(l)}^G \chi_l$ is irreducible; its class is independent of
the choice of $\frak b(l)$; the Kirillov mapping
\begin{gather*}
\kir_G:  \ \ \frak g^* \rightarrow \hat G,\qquad  l \mapsto\pi_l
\end{gather*}
 is surjective and factors to a
bijection $\frak g^*/G \rightarrow \hat G$. Given $\pi \in \hat G$, we write $\Omega(\pi) \in \frak g^*/G$ to denote the inverse
image of $\pi$ under the Kirillov mapping $\kir_G$.


\subsection{Rational structures and uniform subgroups}\label{section2.3}

 In this section we present some results on discrete uniform
subgroups of nilpotent Lie groups.


%{\bf Rational structures.}
 Let $G$ be a nilpotent, connected
and simply connected real Lie group and let $\frak g$ be its Lie
algebra. We say that $\frak g$ (or $G$) has a \textit{rational
structure} if there is a Lie algebra $\frak g_\mathds{Q}$ over
$\mathds{Q}$ such that $\frak g \cong \frak g_\mathds{Q} \otimes
\mathds{R}$. It is clear that $\frak g$ has a rational structure if
and only if $\frak g$ has an $ \mathds{R}$-basis $(X_1,\dots,X_n)$
with rational structure constants.

%{\bf Uniform subgroups.}
 A discrete subgroup $\Gamma$ is
called \textit{uniform} in $G$ if the quotient space $G/\Gamma$ is
compact. The homogeneous space $G/\Gamma$ is called a
\textit{compact nilmanifold}. A proof of the next result can be
found in Theorem~7 of \cite{Malcev1} or in Theorem~2.12 of~\cite{Raghunathan}.


\begin{theorem}[the Malcev rationality criterion] \label{Malcev-cri}
Let $G$ be a simply connected nilpotent Lie
group, and let $\frak g$ be its Lie algebra. Then $G$ admits a
uniform subgroup $\Gamma$ if and only if $\frak g$ admits a basis
$(X_1,\dotsc,X_n)$ such that
\[
[X_i, X_j] =  \sum_{k=1}^n c_{ijk} X_k\qquad \forall\,  1\leq i, j\leq n,
\]
 where the constants $c_{ijk}$ are all
rational.
\end{theorem}



More precisely, if $G$ has a uniform subgroup $\Gamma$, then $\frak
g$  has a rational structure such that
\[
\frak g_\mathds{Q}= \frak g_{\mathds{Q}, \Gamma}= \spannq\set{\log(\Gamma)}.
\]
 Conversely, if
$\frak g$ has a rational structure given by some
$\mathds{Q}$-algebra $\frak g_\mathds{Q}\subset \frak g$, then~$G$
has a uniform subgroup $\Gamma$ such that $\log(\Gamma) \subset
\frak g_\mathds{Q}$ (see \cite{Cor1} or \cite{Malcev1}).





%{\bf Strong Malcev basis.}
\begin{definition}[\cite{Cor1}] Let $\frak g$ be a nilpotent Lie algebra and let
$\bbase=(X_1, \ldots, X_n)$ be a basis of $\frak g$.
 We say that $\bbase$ is a  strong  Malcev
basis for $\frak g$ if  $\frak g_i=\spann\set{X_1, \ldots, X_i}$ is
 an ideal of $\frak g$ for each $1\leq i\leq n$.
\end{definition}


  Let $
\Gamma$ be a discrete uniform subgroup of $G$. A strong Malcev basis
$(X_1,\dotsc,X_n)$ for $\frak g$ is said to be \textit{strongly
based on} $\Gamma$ if
\begin{gather*}%\label{strongly basis}
\Gamma= \exp(\mathds{Z}X_1)\dotsb \exp(\mathds{Z}X_n).
\end{gather*}
 Such a basis always exists (see \cite{Cor1,Matsushima1}).

\subsection{Rational subgroups}\label{2.4}

\begin{definition}[rational subgroup] Let $G$ be a connected
simply connected nilpotent Lie group with Lie algebra~$\frak g$. We
suppose that~$\frak g$ has rational structure given by $\frak
g_\dsq$.
\begin{itemize}\itemsep=0pt
  \item[$(1)$]
Let  $\frak h$ be an $\mathds{R}$-subspace of $\frak g$.  We say
that $\frak h$ is \textit{rational} if $\frak h = \spann\set{\frak
h_\mathds{Q}}$ where  $\frak h_\mathds{Q}= \frak h \cap \frak
g_\mathds{Q}$.
\item[$(2)$]
 A connected, closed subgroup
$H$ of $G$ is called \textit{rational} if its Lie algebra $\frak h$
is rational.
\end{itemize}
\end{definition}


\begin{remark}\label{ratio-intersection}
The $\dsr$-span and the intersection of rational subspaces are
rational \cite[Lemma 5.1.2]{Cor1}.
\end{remark}

\begin{definition}[subgroup with good $\Gamma$-heredity, \cite{Vinberg1}] Let $\Gamma$ be a discrete uniform subgroup in a
locally compact group $G$ and $H$ a closed subgroup in $G$. We say
that $H$ is a subgroup with good $\Gamma$-heredity if the
intersection $\Gamma\cap H$ is a discrete uniform subgroup of $H$.
\end{definition}



\begin{theorem}[\protect{\cite[Lemma A.5]{Cor9}}] Let $G$ be a connected, simply connected
nilpotent Lie group with Lie algebra  $\frak g$,  let $\Gamma$ be a
discrete uniform subgroup of $G$, and give $\frak g$ the rational
structure $\frak g_\dsq = \spannq\set{\log(\Gamma)}$.  Let $H$ be a
Lie subgroup of $G$. Then  the following statements are equivalent
\begin{itemize}\itemsep=0pt
  \item[$(1)$] $H$ is rational;
  \item[$(2)$] $H$ is a subgroup with good $\Gamma$-heredity;
  \item[$(3)$] The group $H$ is $\Gamma$-closed $($i.e., the set
  $H\Gamma$ is closed in $G)$.
\end{itemize}
\end{theorem}




A  proof of the next result can be found in Proposition 5.3.2 of~\cite{Cor1}.

\begin{proposition}\label{suite-adapte-rational}
Let $\Gamma$ be  discrete uniform subgroup in a nilpotent Lie group
$G$, and let $H_1\subsetneqq H_2\subsetneqq\cdots \subsetneqq H_k =G$
be  rational normal subgroups of $G$. Let $\frak h_1, \ldots, \frak
h_{k-1}, \frak h_k =\frak g$ be the corresponding Lie algebras. Then
there exists a strong Malcev basis $(X_1,\ldots, X_n)$ for $\frak g$
strongly based on $\Gamma$ and passing through $\frak h_1, \ldots,
\frak h_{k-1}$.
\end{proposition}


 Now let $G$ be a connected, simply connected nilpotent Lie
group with Lie algebra $\frak g$, and suppose that $\frak g$ has a
rational structure given by  the discrete uniform subgroup $\Gamma$.
A real linear functional $f \in \frak g^*$ is called rational ($f\in
\frak g_\mathds{Q}^*$,  $\frak
g_\mathds{Q}=\spannq\set{\Log{\Gamma}}$) if $\langle f, \frak
g_\mathds{Q}\rangle \subset \mathds{Q}$, or equivalently $\langle f,
\Log{\Gamma} \rangle \subset \mathds{Q}$.


\begin{proposition}[\cite{Cor9}, Theorem A.7] Let $G$ be a nilpotent Lie group with rational
structure and let $\frak g$ be its Lie algebra. If $l\in \frak g^*$
is rational, then its radical $\frak g(l)$ is rational.
\end{proposition}


\section{Preliminary  results}\label{section3}


\begin{definition} A functional $l\in \frak g^*$ is said in general
position or generic linear functional, if its coadjoint orbit has
maximum dimension.
\end{definition}
 The following proposition will be used in the sequel

\begin{proposition}\label{rational-functional}
Let $G = \exp(\frak g)$ be a nilpotent Lie group and $\Gamma$ a
discrete uniform subgroup of~$G$. Then there exist rational generic
linear functionals.
\end{proposition}

\begin{proof}
Let $\mathcal{O}$ be the set of elements in general position in
$\frak g^*$. We have $\mathcal{O}$ is a nonempty Zariski open set
in $\frak g^*$. Since $\frak g_\mathds{Q}^* $ is dense in $\frak
g^*$ then $\frak g_\mathds{Q}^*\cap \mathcal{O} \ne \varnothing$.
\end{proof}


%{\bf Fundamental domains for uniform subgroups.}
 Before stating the next result, we introduce some notations and
def\/initions. We f\/irst recall the concept of fundamental domain.

\begin{definition} Let $G$ be a topological group and let $H$ be a
subgroup of $G$. A fundamental domain for $H$ is a Borel subset
$\Omega$ of $G$ such that
\[
G= \bigsqcup_{h\in H} \Omega h
\] is
the disjoint union of the Borel subsets $ \Omega h$, $h\in H $.
\end{definition}


\begin{remark}
It is clear that a Borel subset  $\Omega$ of $G$ is a
fundamental domain for $H$ in $G$ if and only if  the natural map
$\Omega \rightarrow G/H$ def\/ined by $g\mapsto gH$ is
bijective.
\end{remark}

Let $\Gamma$ be a discrete uniform subgroup of a connected, simply
connected nilpotent Lie group~$G$, and let $ \bbase =(X_1, \ldots,
X_n)$ be a strong Malcev basis for the Lie algebra $\frak g$ of $G$
strongly based on~$\Gamma$. Def\/ine the mapping $\mathbf{E}_{\bbase}:
\mathds{R}^n\rightarrow G$ by
\begin{gather*}
\mathbf{E}_{\bbase}(T) = \exp(t_n X_n)\cdots \exp(t_1 X_1),
\end{gather*}
where $T= (t_1, \ldots, t_n)\in \mathds{R}^n$. It is well known that
 $\mathbf{E}_{\bbase}$ is a dif\/feomorphism~\cite{Cor1}. Let
\begin{gather*}%\label{tore}
\tore = [0, 1) = \set{t\in \mathds{R}:\ 0\leq t < 1}
\end{gather*}
and let
\begin{gather*}%\label{fundamental-domain}
\Omega= \mathbf{E}_{\bbase}(\tore^n).
\end{gather*}
Then $\Omega$ is a fundamental domain for $\Gamma$  in~$G$
\cite[Lemma~3.6]{Cor9}, and the mapping $\mathbf{E}_{\bbase}$ maps
the Lebesgue measure $dt$ on $\tore^n$ to the $G$-invariant
probability measure $\nu$ on $G/\Gamma$, that is, for~$\varphi $ in~$ \mathcal{C}(G/\Gamma)$, we have
\begin{gather*}%\label{measure-fundamental-domain}
\int_{G/\Gamma} \varphi( \dot{g}) d\nu( \dot{g}) = \int_{\tore^n}
\varphi( \mathbf{E}_{\bbase}(t)) dt.
\end{gather*}
Furthermore, for $\phi\in \mathcal{C}_c(G)$ we have
\begin{gather}\label{measure-fundamental-domain1}
\int_{G} \phi(g) d\nu(g) =  \sum_{s\in \dsz^n}\int_{\tore^n}
\phi\big( \mathbf{E}_{\bbase}(t) \mathbf{E}_{\bbase}(s)\big) dt.
\end{gather}


 The following
proposition will be used  in the sequel.

\begin{proposition}\label{decompositionmesure}
Let $G = \exp(\frak g)$ be a connected, simply connected nilpotent
Lie group and let~$\Gamma$ be a discrete uniform subgroup of $G$.
Let  $\frak m$ be a rational ideal of $\frak g$ of dimension $k$.
Let $M= \exp(\frak m)$ and let $(X_1, \ldots, X_n)$ be a strong
Malcev basis of $\frak g$ strongly based on $\Gamma$ passing through
$\frak m $.  For every $\xi\in \mathcal{C}_c(G/M)$, we have
\begin{gather*}%\label{100}
\int_{G/M}  \xi(g)   dg   =
 \sum_{s\in \mathds{Z}^{n-k}}
\int_{\tore^{n-k}}  \xi\Big(\exp(x_{n-k} \overline{X}_n)\cdots
\exp(x_1 \overline{X}_{k+1})\\
\hphantom{\int_{G/M}  \xi(g)   dg   =}{}\times
\exp(s_{n-k} \overline{X}_n)\cdots
\exp(s_1 \overline{X}_{k+1})\Big)
   dx_n\cdots dx_{k+1},
   \end{gather*}
where $\overline{X}_i$ are the image of $X_i$ under the canonical
projection $p: \frak g\rightarrow \frak g/\frak m$.
\end{proposition}

 \begin{proof}
Since $M$ is rational, it follows  from  Lemma 5.1.4 of \cite{Cor1},
that $P(\Gamma)$ is a uniform subgroup of $P(G)$, where $P: G
\rightarrow G/M$ is the canonical projection. On the other hand,
$(\overline{X}_{k+1}, \ldots, \overline{X}_{n})$ is a strong Malcev
basis of $\frak g/\frak m$ strongly based on $P(\Gamma)$. We
conclude by applying (\ref{measure-fundamental-domain1}).
 \end{proof}



\section{On the rational structure  of
one-parameter\\ metabelian nilpotent Lie groups}\label{section4}


\begin{definition}[one-parameter metabelian nilpotent Lie algebra]
A nonabelian nilpotent Lie algebra  $\frak g$ is said to be one
parameter metabelian nilpotent Lie algebra, if it admits  a
co-dimensional one abelian ideal
 in $\frak g$. \end{definition}


First, we give an important example of one-parameter metabelian
nilpotent Lie groups.

 \begin{example}[the generic f\/iliform nilpotent Lie group]
  Let $G$ be the generic f\/iliform
nilpotent Lie group of dimension $n$ with Lie algebra $\frak g$,
where
\begin{gather*}
 \frak g = \spann\set{X_1,
\ldots, X_n}
\end{gather*}
The Lie brackets  given by
\begin{gather*}
[X_n, X_i] = X_{i-1},  \qquad i=2,\ldots,  n-1,
\end{gather*}
and the nondef\/ined brackets being equal to zero or obtained by
antisymmetry. It is clear that~$\frak g$ is a one-parameter
metabelian nilpotent Lie algebra.
\end{example}

\begin{remark} Let $\frak g$ be a one-parameter metabelian
 nilpotent Lie algebra. Any
codimension one abelian ideal $\frak m\subset \frak g$ is a common
polarization for all functionals $l\in \frak g^*$ in general
position (i.e.,  $l\vert_{[\frak g, \frak g]}\ne 0$).
\end{remark}


\begin{definition}[one-parameter metabelian nilmanifold]\sloppy
A factor space of a one-parameter met\-abe\-lian nilpotent Lie group
over a discrete uniform subgroup is called a one-parameter
metabelian nilmanifold.
\end{definition}


\begin{proposition}\label{decomposition-speciale} Let $\frak g$ be a
one-parameter metabelian nilpotent Lie algebra. Then we have the
following decomposition
\begin{gather*}
\frak g = \mathds{R} X \oplus \oplus_{i=1}^p \mathcal{L}_{n_i}
\oplus \frak a
\end{gather*}
such that for all $i=1,\ldots, p$, the subalgebra $\mathds{R} X
\oplus  \mathcal{L}_{n_i} $ is the generic filiform nilpotent Lie
algebra of dimension $n_i+1$ and $\frak a \subset \frak z(\frak g)$.
 \end{proposition}

 \begin{proof} Let $\frak I$ be a one-codimensional abelian ideal of
 $\frak  g$ and let $X\in \frak g$ such that
\[
\frak g= \frak I\oplus \spann\set{X}.
\]
 Let $\mathrm{ad} X\vert_{\frak I}$ be the
 restriction of $\mathrm{ad}
 X$ to $\frak I$;
\[
\mathrm{ad}\, X\vert_{\frak I}: \quad \frak I\rightarrow \frak I,\qquad
 Y\mapsto [X, Y].
 \]
 Remark that $\mathrm{ad}\, X$ acts as a nilpotent linear
 transformation on~$\frak I$.
 By the Jordan normal form theorem, the matrix of
 $\mathrm{ad} X\vert_{\frak I}$ is similar to a matrix in real Jordan canonical form. Then
 there exist $\mathscr{B}_{\frak I}=(e_1, \ldots, e_{n-1})$  a
 basis of $\frak I$, $J_{0, n_i}$, $1\leq i\leq s$  elementary Jordan
 blocks of order $n_i$ such that
\[
\mathrm{Mat}(\mathrm{ad}\, X\vert_{\frak
 f}, \mathscr{B}_{\frak I}) =
  \mathrm{diag}[J_{0, n_1}, \ldots, J_{0, n_s}].
\]
  This completes the proof.
 \end{proof}

\begin{proposition}
Let $G$ be a one-parameter metabelian nilpotent Lie group with Lie
algebra~$\frak g$.
 Then $G$ admits a discrete uniform subgroup.
\end{proposition}

\begin{proof} This follows at once from
 the  Malcev  rationality criterion
 (Theorem~\ref{Malcev-cri}) and Proposition~\ref{decomposition-speciale}.
 \end{proof}

\begin{proposition}\label{maryem10}
Let $G$ be a one-parameter metabelian nilpotent Lie group of
dimension $n$.
 Let~$\Gamma $ be a discrete uniform subgroup of $G$. Then $G$ admits a
rational
 abelian ideal of codimension one.
\end{proposition}

Before proving the proposition, we need the following theorem.
\begin{theorem}[\protect{\cite[p.~17]{Bernat}}]
If $l\in \frak g^*$ is in general position then $\frak g(l)$ is
abelian.
\end{theorem}

\begin{proof}[Proof of Proposition~\ref{maryem10}]
Let $(e_1, \ldots, e_{n})$ be a strong Malcev basis for $\frak g$
strongly based on~$\Gamma$. Let $l$ be a rational  linear functional
in general position (see Proposition \ref{rational-functional}).
 The stabilizer of $l$ is a~rational abelian  subalgebra. On the other hand, applying equality~(\ref{dim-pilarisation}) we obtain
\begin{gather}\label{dim-stab}
\dim{\frak g(l)}=
 n-2.
 \end{gather}
  Let $(e_{i_1}, e_{i_2})$  be a basis of
$\frak g$ modulo $\frak g(l)$. If $ [ e_{i_1}, \frak g(l)] = \lbrace
0\rbrace$, then $\frak g(l)\oplus \spann\lbrace e_{i_1}\rbrace$  is
a rational abelian ideal of $\frak g$ of codimension one. If $ [
e_{i_1}, \frak g(l)] \ne \lbrace 0 \rbrace$, let
\[
\frak a = \lbrace X \in \spann\lbrace e_{i_1}, e_{i_2}\rbrace:\
[ X, \frak g(l)] =\set{0}\rbrace.
\] In this case, we have $\dim
{\frak a} =1$. Moreover, it is clear that
\[
 \frak a = \lbrace X \in \frak g:\
[ X, \frak g(l)] =\set{0}\rbrace\cap \spann\lbrace e_{i_1},
e_{i_2}\rbrace.
\]
 By Proposition~5 of~\cite{Ghorbel1}, the space
\[
\lbrace X \in \frak g:\
[ X, \frak g(l)] =\set{0}\rbrace
\] is rational. Consequently, $\frak
a$ is rational subspace in~$\frak g$. Thus $\frak a \oplus \frak
g(l)$ is a rational abelian ideal of~$\frak g$ of codimension one.
\end{proof}
